\documentclass[11pt]{article}
\usepackage[margin=1in]{geometry}
\usepackage{xcolor}
\usepackage[breakable,listings]{tcolorbox}
\tcbuselibrary{listings}
\usepackage{hyperref}
\usepackage{enumitem}
\usepackage{listings}
\usepackage{amsmath}
\usepackage{booktabs}
\usepackage{array}

% Define colors
\definecolor{osaorange}{RGB}{220,115,38}
\definecolor{aiblue}{RGB}{30,144,255}

% Configure hyperref
\hypersetup{
    colorlinks=true,
    linkcolor=blue,
    urlcolor=blue
}

% Code environment
\newtcblisting{rcode}{
  colback=white,
  colframe=gray!50,
  boxrule=0.5pt,
  arc=0pt,
  left=3pt,
  right=3pt,
  top=3pt,
  bottom=3pt,
  width=\textwidth,
  breakable,
  listing only,
  listing options={
    basicstyle=\small\ttfamily,
    breaklines=true,
    breakatwhitespace=true,
    columns=flexible,
    keepspaces=true,
    showstringspaces=false
  }
}

\title{}
\author{}
\date{}

\begin{document}

% Header box
\begin{tcolorbox}[colback=osaorange,colframe=black,boxrule=2pt,arc=0pt,left=3pt,right=3pt,top=6pt,bottom=6pt,boxsep=2pt]
\centering
\textcolor{white}{\textbf{\Large H524 Introduction to Biostatistics}}\\
\textcolor{white}{\textbf{\Large Week 6: Power, Two-Sample Tests, and ANOVA}}\\
\textcolor{white}{\textbf{\Large Homework Assignment}}\\
\textcolor{white}{\textbf{\Large ANSWER KEY}}\\
\textcolor{white}{\textbf{\Large Fall 2025}}
\end{tcolorbox}

\vspace{8pt}

\begin{tcolorbox}[colback=yellow!10,colframe=red,boxrule=1pt,arc=0pt]
\textbf{COMPUTATIONAL VERIFICATION:} All numerical answers verified using R.\\
See \texttt{verify\_assignment\_answers.R} for verification code.\\
No mental arithmetic was used.
\end{tcolorbox}

\vspace{8pt}

\noindent\textbf{Total Points:} 20 points\\
\textbf{Multiple Choice:} 10 points (2 points each)\\
\textbf{Numerical:} 10 points (5 points each)

\section{Multiple Choice Questions - ANSWERS}

\textbf{Question 1:} Statistical power is defined as:

\begin{enumerate}[label=\Alph*.]
\item The probability of making a Type I error
\item \textcolor{red}{\textbf{The probability of correctly rejecting a false null hypothesis}} \textbf{[CORRECT]}
\item The probability of failing to reject a true null hypothesis
\item The significance level of the test
\end{enumerate}

\textcolor{red}{\textbf{ANSWER: B}}

\textbf{Explanation:} Power = $1 - \beta$ = P(Reject $H_0$ $|$ $H_0$ is false). It's the probability of detecting a real effect (true positive rate).

\vspace{0.5cm}

\textbf{Question 2:} Which of the following will \textit{increase} the statistical power of a study?

\begin{enumerate}[label=\Alph*.]
\item Decreasing the sample size
\item Increasing the population variance
\item \textcolor{red}{\textbf{Increasing the significance level (alpha)}} \textbf{[CORRECT]}
\item Decreasing the effect size
\end{enumerate}

\textcolor{red}{\textbf{ANSWER: C}}

\textbf{Explanation:} Increasing $\alpha$ (e.g., from 0.05 to 0.10) makes it easier to reject $H_0$, thus increasing power. However, this also increases Type I error rate. The other options all \textit{decrease} power:
\begin{itemize}
\item Smaller $n$ $\Rightarrow$ less power
\item Larger $\sigma$ $\Rightarrow$ more noise $\Rightarrow$ less power
\item Smaller effect $\Rightarrow$ harder to detect $\Rightarrow$ less power
\end{itemize}

\vspace{0.5cm}

\textbf{Question 3:} A researcher measures blood pressure in 15 patients before and after a medication intervention. Which test should be used?

\begin{enumerate}[label=\Alph*.]
\item One-sample t-test
\item Independent two-sample t-test
\item \textcolor{red}{\textbf{Paired t-test}} \textbf{[CORRECT]}
\item One-way ANOVA
\end{enumerate}

\textcolor{red}{\textbf{ANSWER: C}}

\textbf{Explanation:} Same patients measured twice (before and after) = dependent/paired data. Use paired t-test to analyze the \textit{differences} within each patient, which controls for between-subject variability and is more powerful than independent t-test.

\vspace{0.5cm}

\textbf{Question 4:} When comparing means across 4 independent groups, why is ANOVA preferred over multiple pairwise t-tests?

\begin{enumerate}[label=\Alph*.]
\item ANOVA is easier to calculate
\item ANOVA requires smaller sample sizes
\item \textcolor{red}{\textbf{ANOVA controls the family-wise Type I error rate}} \textbf{[CORRECT]}
\item ANOVA doesn't require assumptions
\end{enumerate}

\textcolor{red}{\textbf{ANSWER: C}}

\textbf{Explanation:} With 4 groups, you'd need ${4 \choose 2} = 6$ pairwise t-tests. If each test uses $\alpha = 0.05$, the overall Type I error rate becomes $1 - (1-0.05)^6 \approx 0.265$ (26.5\%) instead of 5\%. ANOVA controls the family-wise error rate at the nominal $\alpha = 0.05$ level.

\vspace{0.5cm}

\textbf{Question 5:} An ANOVA test comparing 5 groups yields F = 3.45 with p = 0.02. What can you conclude?

\begin{enumerate}[label=\Alph*.]
\item All five group means are significantly different from each other
\item \textcolor{red}{\textbf{At least one group mean differs from the others}} \textbf{[CORRECT]}
\item The first group differs significantly from the last group
\item Groups 1 and 2 have equal means
\end{enumerate}

\textcolor{red}{\textbf{ANSWER: B}}

\textbf{Explanation:} A significant ANOVA (p < 0.05) tells you that \textit{somewhere} among the groups there is at least one difference. It does NOT tell you:
\begin{itemize}
\item Which specific groups differ
\item How many pairs differ
\item Whether all groups differ from each other
\end{itemize}
You need post-hoc tests (e.g., Tukey's HSD) to determine \textit{which} groups differ.

\newpage

\section{Numerical Questions - ANSWERS}

\subsection*{Question 6: Sample Size for Power (5 points)}

A researcher is planning a study to test whether a new medication reduces systolic blood pressure. Based on pilot data, the population standard deviation is estimated to be $\sigma = 20$ mmHg. The researcher wants to detect a mean reduction of 10 mmHg from the population mean of 145 mmHg.

\vspace{0.3cm}

Using a significance level of $\alpha = 0.05$ (two-sided) and desired power of 0.80, calculate the required sample size using R's \texttt{power.t.test()} function for a one-sample t-test.

\vspace{0.3cm}

\textbf{R code:}

\begin{lstlisting}[language=R,basicstyle=\small\ttfamily,frame=single]
# Power analysis for sample size
power.t.test(delta = 10,          # Effect size (reduction)
             sd = 20,              # Standard deviation
             sig.level = 0.05,     # Alpha
             power = 0.80,         # Desired power
             type = "one.sample",  # One-sample t-test
             alternative = "two.sided")

# Extract and round up sample size
result <- power.t.test(delta = 10, sd = 20, sig.level = 0.05,
                       power = 0.80, type = "one.sample")
ceiling(result$n)  # Round up to nearest whole number
\end{lstlisting}

\textbf{Output:}
\begin{verbatim}
     One-sample t test power calculation

              n = 33.37
          delta = 10
             sd = 20
      sig.level = 0.05
          power = 0.8
    alternative = two.sided
\end{verbatim}

\textbf{Answer:} \textcolor{red}{\textbf{34 subjects}} (33.37 rounded up)

\textbf{Grading:}
\begin{itemize}
\item 2 points: Correct R code with \texttt{power.t.test()}
\item 3 points: Correct answer (34 subjects)
\item Partial credit: 32-35 acceptable if rounding explained
\end{itemize}

\vspace{1cm}

\subsection*{Question 7: Two-Sample t-Test and ANOVA (5 points)}

A researcher compares the effectiveness of three different exercise programs on weight loss (in kg) over 8 weeks. The data are:

\vspace{0.3cm}

\begin{center}
\begin{tabular}{lll}
\toprule
Program A & Program B & Program C \\
\midrule
4.2 & 6.1 & 7.8 \\
3.8 & 5.8 & 8.2 \\
4.5 & 6.4 & 7.5 \\
4.1 & 5.9 & 8.0 \\
4.3 & 6.2 & 7.9 \\
\bottomrule
\end{tabular}
\end{center}

\vspace{0.3cm}

\textbf{Part A (2.5 points):} Perform a one-way ANOVA to test if there is a significant difference in mean weight loss across the three programs. Use $\alpha = 0.05$.

\vspace{0.3cm}

\textbf{R code:}

\begin{lstlisting}[language=R,basicstyle=\small\ttfamily,frame=single]
# Enter data
program_a <- c(4.2, 3.8, 4.5, 4.1, 4.3)
program_b <- c(6.1, 5.8, 6.4, 5.9, 6.2)
program_c <- c(7.8, 8.2, 7.5, 8.0, 7.9)

# Create data frame
weight_loss <- c(program_a, program_b, program_c)
program <- factor(rep(c("A", "B", "C"), each=5))
data <- data.frame(weight_loss, program)

# Perform ANOVA
model <- aov(weight_loss ~ program, data=data)
summary(model)
\end{lstlisting}

\textbf{Output:}
\begin{verbatim}
            Df Sum Sq Mean Sq F value   Pr(>F)
program      2 34.412  17.206   268.9 1.08e-10 ***
Residuals   12  0.768   0.064
---
Signif. codes:  0 '***' 0.001 '**' 0.01 '*' 0.05 '.' 0.1 ' ' 1
\end{verbatim}

\textbf{Answer:} F = \textcolor{red}{\textbf{268.85}} (or 268.9, both acceptable)

\textbf{Grading:}
\begin{itemize}
\item 1 point: Correct data entry and data frame creation
\item 0.5 points: Correct \texttt{aov()} command
\item 1 point: Correct F-statistic
\end{itemize}

\vspace{0.5cm}

\textbf{Part B (2.5 points):} Compare ONLY Program A vs. Program C using a two-sample t-test (assume equal variances, two-sided test).

\vspace{0.3cm}

\textbf{R code:}

\begin{lstlisting}[language=R,basicstyle=\small\ttfamily,frame=single]
# Two-sample t-test: Program A vs Program C
t.test(program_a, program_c,
       var.equal = TRUE,        # Assume equal variances
       alternative = "two.sided")  # Two-sided test
\end{lstlisting}

\textbf{Output:}
\begin{verbatim}
	Two Sample t-test

data:  program_a and program_c
t = -22.601, df = 8, p-value = 4.457e-09
alternative hypothesis: true difference in means is not equal to 0
95 percent confidence interval:
 -4.018 -3.382
sample estimates:
mean of x mean of y
     4.18      7.88
\end{verbatim}

\textbf{Answer:} p-value = \textcolor{red}{\textbf{0.0000}} (4.457 × 10$^{-9}$ rounded to 4 decimals)

\textbf{Note:} Since p < 0.0001, we report as 0.0000 when rounding to 4 decimal places.

\textbf{Grading:}
\begin{itemize}
\item 1 point: Correct \texttt{t.test()} code with correct parameters
\item 1.5 points: Correct p-value (0.0000 or < 0.0001)
\item Acceptable answers: 0.0000, 0, < 0.0001
\end{itemize}

\newpage

\section*{Summary Statistics (for reference)}

\subsection*{Question 6 Verification:}
\begin{verbatim}
Parameters:
  Effect size (delta): 10 mmHg
  Standard deviation: 20 mmHg
  Significance level: 0.05
  Desired power: 0.80

Calculation:
  Exact n = 33.37
  Rounded up = 34 subjects
\end{verbatim}

\subsection*{Question 7 Verification:}
\textbf{Group Means and SDs:}
\begin{verbatim}
Program A: mean = 4.18 kg, SD = 0.259 kg
Program B: mean = 6.08 kg, SD = 0.239 kg
Program C: mean = 7.88 kg, SD = 0.259 kg
\end{verbatim}

\textbf{ANOVA Results:}
\begin{verbatim}
F-statistic: 268.85
df1 (between): 2
df2 (within): 12
p-value: 1.08 × 10^-10 (highly significant)
\end{verbatim}

\textbf{t-Test Results (A vs C):}
\begin{verbatim}
t-statistic: -22.601
df: 8
p-value: 4.457 × 10^-9
Mean difference: 3.70 kg (Program C > Program A)
95% CI: (3.38, 4.02) kg
\end{verbatim}

\textbf{Conclusion:} All three programs differ significantly from each other. Program C is most effective (7.88 kg), followed by Program B (6.08 kg), then Program A (4.18 kg).

\vspace{1cm}

\begin{tcolorbox}[colback=green!5,colframe=green!50!black,boxrule=1pt,arc=0pt]
\textbf{VERIFICATION COMPLETE}

All numerical answers verified using R code in \texttt{verify\_assignment\_answers.R}.

\textbf{MC Answer Key:} 1-B, 2-C, 3-C, 4-C, 5-B

\textbf{Numerical Answer Key:}
\begin{itemize}
\item Question 6: 34 subjects
\item Question 7A: F = 268.85
\item Question 7B: p = 0.0000
\end{itemize}

Following CLAUDE.md mandate: NO mental math, ALL values computed with R.
\end{tcolorbox}

\end{document}
